%% ============================================================
%%  Aula 9 — Aprendizagem por reforço eficiente em dados I:
%%           bandidos multi-braços, arrependimento e otimismo
%%  Adaptação em pt-BR do CS234 (Stanford, Emma Brunskill), Lecture 9
%%  Prof. Marcelo Keese Albertini — Universidade Federal de Uberlândia
%%  Compilar com XeLaTeX (2x) — latexmk -xelatex aula09.tex
%% ============================================================
\documentclass[aspectratio=169, 10pt]{beamer}

\usetheme{AprendizadoRL}      % ANTES do babel — fontspec precisa vir primeiro
\usepackage{babel}
\babelprovide[import, main]{brazilian}
\usepackage{booktabs}
\usepackage{pgfplots}
\pgfplotsset{compat=1.18}

\title[Aula 9 — Bandidos e arrependimento]{Aula 9: Aprendizagem por reforço\\ eficiente em dados I}
\subtitle[Aprendizagem por Reforço]{Bandidos multi-braços \textbullet\ Arrependimento \textbullet\ $\epsilon$-guloso \textbullet\ Otimismo e UCB}
\author[Prof. Marcelo K. Albertini]{Prof. Marcelo Keese Albertini \\
  \footnotesize Universidade Federal de Uberlândia \\
  \footnotesize adaptado de Emma Brunskill, CS234 (Stanford)}
\institute{Universidade Federal de Uberlândia}
\date{2026}

% ---- Nota discreta: as perguntas ficam para o professor conduzir em aula
\newcommand{\paradiscussao}{\smallskip
  {\footnotesize\color{rlCinza}\itshape Pergunta para discussão conduzida em aula.}}

% ---- Macros locais desta aula
\newcommand{\Qhat}{\widehat{Q}}
\newcommand{\muhat}{\hat{\mu}}
\newcommand{\Nc}{N}                         % contagem
\newcommand{\gap}{\Delta}
\newcommand{\UCB}{\mathrm{UCB}}
\newcommand{\LCB}{\mathrm{LCB}}
\newcommand{\KL}{D_{\mathrm{KL}}}
\newcommand{\one}{\mathbb{1}}
\newcommand{\Reg}{L}
\newcommand{\Ber}{\mathrm{Bernoulli}}
\newcommand{\Bet}{\mathrm{Beta}}

% ---- Estilos TikZ locais
\tikzset{
  caixatit/.style={draw=rlAzul!35, very thick, rounded corners=2.5pt,
                   fill=rlAzul!5, align=center, inner sep=3mm,
                   text width=27mm, font=\small},
  braco/.style={draw=rlAzul, very thick, fill=rlAzul!8, rounded corners=2pt,
                minimum width=9mm, minimum height=5mm, font=\scriptsize},
}

\begin{document}

\begin{frame}[plain, noframenumbering]\titlepage\end{frame}

%% ============================================================
\begin{frame}{Retomando a aula anterior}
  \framesubtitle{Aprendizagem a partir de preferências humanas}

  \begin{quizbox}[fontupper=\footnotesize]{}
    Quais afirmações são verdadeiras?
    \begin{enumerate}\setlength{\itemsep}{0ex}
      \item RLHF e DPO ambos aprendem uma representação \alert{explícita}
            de um modelo de recompensa a partir de preferências.
      \item Ambos estão limitados a ser, no máximo, tão bons quanto os
            melhores exemplos dos pares de preferência.
      \item O DPO não utiliza política de referência.
    \end{enumerate}
    \paradiscussao
  \end{quizbox}

  \pause

  \begin{respbox}[fontupper=\footnotesize]
    \textbf{Nenhuma das três.} (1) O DPO dispensa o modelo explícito: a
    recompensa fica \emph{implícita} em
    $\log\frac{\pol_\theta}{\pol_{\text{ref}}}$. (2) Otimizar uma recompensa
    aprendida permite \emph{superar} os exemplos vistos. (3)
    $\pol_{\text{ref}}$ é parte constitutiva da perda do DPO.
  \end{respbox}
\end{frame}

%% ============================================================
\begin{frame}{Onde estamos no curso}
  \begin{columns}[T]
    \begin{column}{0.52\textwidth}
      \begin{itemize}
        \item \textbf{Aulas 1--2:} planejamento com modelo conhecido.
        \item \textbf{Aulas 3--4:} avaliação e controle sem modelo.
        \item \textbf{Aulas 5--7:} gradientes de política, métodos modernos.
        \item \textbf{Aula 8:} aprendizagem a partir de preferências.
        \item \textbf{Aulas 9--10:} \alert{eficiência de dados} — quantas
              interações são \emph{necessárias}?
      \end{itemize}
    \end{column}
    \begin{column}{0.46\textwidth}
      \begin{ideiabox}
        Até agora perguntamos \emph{``o algoritmo converge para a política
        ótima?''}. A partir de hoje perguntamos algo mais exigente:
        \emph{``quantos erros ele comete no caminho?''}
      \end{ideiabox}
    \end{column}
  \end{columns}

  \bigskip
  \begin{alertabox}
    Convergência assintótica é um critério barato. Um algoritmo que testa
    todas as ações um milhão de vezes cada e depois age otimamente
    ``converge'' — e é inútil se cada teste custa um paciente, um cliente
    ou uma hora de robô.
  \end{alertabox}
\end{frame}

%% ============================================================
\begin{frame}{Os quatro ingredientes, revisitados}
  \begin{columns}[T]
    \begin{column}{0.5\textwidth}
      A aprendizagem por reforço envolve:
      \begin{itemize}
        \item \textbf{Otimização} — encontrar boas decisões.
        \item \textbf{Generalização} — inferir sobre o que não foi visto.
        \item \textbf{Consequências adiadas} — hoje afeta amanhã.
        \item \alert{\textbf{Exploração}} — decidir o que experimentar.
      \end{itemize}
    \end{column}
    \begin{column}{0.48\textwidth}
      \begin{defbox}{Estratégia desta aula}
        Isolar a \textbf{exploração} removendo temporariamente as
        consequências adiadas. O objeto resultante — o \emph{bandido
        multi-braços} — é simples o bastante para admitir análise exata,
        e as ideias sobrevivem intactas quando as consequências adiadas
        voltam (Aula 10).
      \end{defbox}
    \end{column}
  \end{columns}
\end{frame}

%% ============================================================
\begin{frame}{O que vamos construir hoje}
  \begin{columns}[T]
    \begin{column}{0.33\textwidth}
      {\color{rlTeal}\bfseries CENÁRIOS}\\[0.3em]
      {\small Onde o agente age.}
      \begin{itemize}\small
        \item Bandidos (decisão única)
        \item MDPs (Aula 10)
      \end{itemize}
    \end{column}
    \begin{column}{0.33\textwidth}
      {\color{rlAzulClaro}\bfseries ARCABOUÇOS}\\[0.3em]
      {\small Como se mede qualidade.}
      \begin{itemize}\small
        \item Arrependimento
        \item Arrependimento bayesiano
        \item PAC (Aula 10)
      \end{itemize}
    \end{column}
    \begin{column}{0.33\textwidth}
      {\color{rlLaranja}\bfseries ABORDAGENS}\\[0.3em]
      {\small Classes de algoritmos.}
      \begin{itemize}\small
        \item Guloso, $\epsilon$-guloso
        \item Otimismo (UCB)
        \item Casamento de probabilidade
      \end{itemize}
    \end{column}
  \end{columns}

  \bigskip
  \begin{ideiabox}
    A matriz cenário $\times$ arcabouço $\times$ abordagem não é preenchida
    ponto a ponto: \textbf{uma mesma abordagem costuma atender vários
    arcabouços em vários cenários}. O otimismo, que veremos hoje em
    bandidos, reaparecerá quase sem alteração em MDPs.
  \end{ideiabox}
\end{frame}

%% ============================================================
\section{Bandidos multi-braços}

\begin{frame}{Como julgar um algoritmo de aprendizagem}
  Critérios que já usamos, em ordem crescente de exigência:
  \begin{enumerate}
    \item \textbf{Custo computacional} por iteração.
    \item \textbf{Converge?} Existe um ponto fixo e o algoritmo o alcança?
    \item \textbf{Converge para o \emph{ótimo}?}
    \item \textbf{Com que rapidez} alcança a política ótima?
    \item \alert{\textbf{Quantos erros comete no caminho?}}
  \end{enumerate}

  \bigskip
  \begin{alertabox}
    O critério (5) é o único que fala do \emph{custo da própria
    aprendizagem}. Em ensaios clínicos, precificação, recomendação e
    robótica, ele é o que importa: cada decisão subótima é um dano real,
    não um número numa curva de treino.
  \end{alertabox}
\end{frame}

%% ============================================================
\begin{frame}{Bandido multi-braços}
  \begin{defbox}{Bandido multi-braços}
    Uma tupla $(\gA, \mathcal{R})$ onde
    \begin{itemize}\setlength{\itemsep}{0.2ex}
      \item $\gA$ é um conjunto \emph{conhecido} de $m$ ações (``braços'');
      \item $\mathcal{R}_a(r) = \Prob[r \mid a]$ é uma distribuição de
            recompensas \emph{desconhecida}, uma para cada braço.
    \end{itemize}
    A cada passo o agente escolhe $a_t \in \gA$; o ambiente devolve
    $r_t \sim \mathcal{R}_{a_t}$. Objetivo: maximizar $\sum_{\tau} r_\tau$.
  \end{defbox}

  \begin{columns}[T]
    \begin{column}{0.54\textwidth}
      \begin{itemize}\small
        \item Não há estado: a ação de hoje não altera o problema de amanhã.
        \item O único acoplamento entre passos é \textbf{informacional}:
              o que aprendi ontem muda o que escolho hoje.
      \end{itemize}
    \end{column}
    \begin{column}{0.44\textwidth}
      \centering
      \begin{tikzpicture}[scale=0.78, transform shape]
        \node[estado destacado, minimum size=11mm] (s) {$s$};
        \foreach \i/\ang/\col in {1/140/rlLaranja, 2/180/rlAzulClaro, 3/220/rlTeal}{
          \draw[trans, \col, very thick] (s) to[out=\ang, in=\ang-60, looseness=9] node[pos=0.5, left, font=\scriptsize, text=\col] {$a_\i$} (s);}
        \node[font=\scriptsize, text=rlCinza, below=7mm of s, align=center]
          {MDP de \textbf{um único estado}:\\ todas as ações voltam a $s$.};
      \end{tikzpicture}
    \end{column}
  \end{columns}
\end{frame}

%% ============================================================
\begin{frame}{Bandido, bandido contextual, MDP}
  \centering
  \begin{tikzpicture}[node distance=6mm]
    \node[caixatit] (b) {\textbf{Bandido}\\[0.2em]
      \scriptsize um estado\\ \scriptsize sem consequência adiada\\[0.3em]
      \scriptsize $a \to r$};
    \node[caixatit, right=14mm of b] (cb) {\textbf{Bandido contextual}\\[0.2em]
      \scriptsize estado sorteado i.i.d.\\ \scriptsize sem consequência adiada\\[0.3em]
      \scriptsize $s \to a \to r$};
    \node[caixatit, right=14mm of cb] (mdp) {\textbf{MDP}\\[0.2em]
      \scriptsize estado evolui\\ \scriptsize consequência adiada\\[0.3em]
      \scriptsize $s \to a \to r,\, s'$};
    \draw[trans destac] (b) -- node[above, font=\scriptsize, text=rlCinza] {+ contexto} (cb);
    \draw[trans destac] (cb) -- node[above, font=\scriptsize, text=rlCinza] {+ dinâmica} (mdp);
  \end{tikzpicture}

  \vspace{1em}
  \begin{columns}[T]
    \begin{column}{0.48\textwidth}
      \begin{ideiabox}
        A exploração é \emph{difícil} nos três casos, mas por razões
        diferentes. No bandido, basta decidir \emph{o que testar}. No MDP,
        é preciso decidir também \emph{como chegar} ao lugar onde vale a
        pena testar --- exploração profunda.
      \end{ideiabox}
    \end{column}
    \begin{column}{0.48\textwidth}
      \begin{itemize}
        \item Anúncios, notícias e recomendação: bandidos contextuais, na
              prática, movimentam bilhões.
        \item Ensaios clínicos adaptativos: bandidos puros com
              restrições éticas.
        \item Toda a teoria de hoje é o caso base das cotas de amostra
              para MDPs.
      \end{itemize}
    \end{column}
  \end{columns}
\end{frame}

%% ============================================================
\begin{frame}{Exemplo condutor: como tratar um dedo do pé quebrado}
  \begin{columns}[T]
    \begin{column}{0.55\textwidth}
      \begin{itemize}
        \item Decidir o melhor tratamento para pacientes com fratura de
              dedo do pé.
        \item Três opções: \textbf{(1)} cirurgia, \textbf{(2)}
              imobilização com esparadrapo no dedo vizinho
              (\emph{buddy taping}), \textbf{(3)} não fazer nada.
        \item Desfecho binário: o dedo consolidou (\;$+1$\;) ou não
              (\;$0$\;) após seis semanas, avaliado por radiografia.
        \item Cada braço é uma variável de Bernoulli com parâmetro
              desconhecido $\theta_i$.
      \end{itemize}
    \end{column}
    \begin{column}{0.43\textwidth}
      \begin{alertabox}
        Exemplo \textbf{fictício}. Os números abaixo não representam a
        eficácia real de nenhum tratamento ortopédico.
      \end{alertabox}
      \begin{defbox}{Parâmetros verdadeiros (desconhecidos)}
        \centering\small
        \begin{tabular}{lc}
          cirurgia $a_1$     & $\theta_1 = 0{,}95$ \\
          esparadrapo $a_2$  & $\theta_2 = 0{,}90$ \\
          nada $a_3$         & $\theta_3 = 0{,}10$ \\
        \end{tabular}
      \end{defbox}
    \end{column}
  \end{columns}
\end{frame}

%% ============================================================
\begin{frame}{Teste sua compreensão: modelando o problema}
  \begin{quizbox}[fontupper=\footnotesize]{}
    Modelando o tratamento como bandido de 3 braços, cada um Bernoulli
    com parâmetro desconhecido $\theta_i$, quais são verdadeiras?
    \begin{enumerate}\setlength{\itemsep}{0ex}
      \item Puxar um braço / tomar uma ação corresponde a o dedo ter
            consolidado ou não.
      \item Um bandido descreve melhor este problema do que um MDP porque
            tratar cada paciente envolve múltiplas decisões.
      \item Depois de tratar um paciente, se $\theta_i \ne 0$ e
            $\theta_i \ne 1$ para todo $i$, às vezes o dedo consolida e às
            vezes não.
    \end{enumerate}
    \paradiscussao
  \end{quizbox}

  \pause

  \begin{respbox}[fontupper=\footnotesize]
    Apenas a \textbf{(3)}. \;
    \textbf{(1)} é falsa: puxar o braço é \emph{escolher o tratamento} --- a
    consolidação é a \emph{recompensa}. \;
    \textbf{(2)} tem a conclusão certa pelo motivo errado: o bandido serve
    porque cada paciente envolve \emph{uma única} decisão. \;
    \textbf{(3)} é a definição de recompensa estocástica.
  \end{respbox}
\end{frame}

%% ============================================================
\begin{frame}{O algoritmo guloso}
  \begin{algbox}{Guloso}
    Estime o valor de cada ação por avaliação de Monte Carlo:
    \[
      \Qhat_t(a) \;=\; \frac{1}{\Nc_t(a)} \sum_{i=1}^{t} r_i\,\one(a_i = a)
      \;\approx\; Q(a) = \E[R(a)],
    \]
    e escolha sempre a ação de maior valor estimado:
    \[
      a_t \;=\; \argmax_{a \in \gA} \Qhat_{t-1}(a).
    \]
  \end{algbox}

  \begin{columns}[T]
    \begin{column}{0.55\textwidth}
      Forma incremental, sem guardar o histórico:
      \[
        \Qhat_{t}(a_t) \;=\; \Qhat_{t-1}(a_t)
          + \frac{1}{\Nc_t(a_t)}\bigl(r_t - \Qhat_{t-1}(a_t)\bigr).
      \]
    \end{column}
    \begin{column}{0.43\textwidth}
      \begin{ideiabox}
        O mesmo molde das aulas anteriores, com $\alpha = 1/\Nc$: a média
        exata.
      \end{ideiabox}
    \end{column}
  \end{columns}
\end{frame}

%% ============================================================
\begin{frame}{O guloso trava}
  \framesubtitle{Uma rodada de amostragem inicial, e o destino selado}

  \begin{columns}[T]
    \begin{column}{0.52\textwidth}
      \begin{algbox}{Execução}
        \textbf{1.} Amostre cada braço uma vez:
        \begin{itemize}\setlength{\itemsep}{0.1ex}\small
          \item $a_1$, $r \sim \Ber(0{,}95)$, obtém $\mathbf{0}$
                $\Rightarrow \Qhat(a_1) = 0$
          \item $a_2$, $r \sim \Ber(0{,}90)$, obtém $\mathbf{+1}$
                $\Rightarrow \Qhat(a_2) = 1$
          \item $a_3$, $r \sim \Ber(0{,}10)$, obtém $\mathbf{0}$
                $\Rightarrow \Qhat(a_3) = 0$
        \end{itemize}
        \textbf{2.} Qual a probabilidade de o guloso escolher cada braço
        em seguida? (empates resolvidos uniformemente)
      \end{algbox}
      \pause
      \begin{respbox}
        $\Prob(a_2) = 1$. O máximo é único.
      \end{respbox}
    \end{column}
    \begin{column}{0.46\textwidth}
      \pause
      \begin{alertabox}
        \textbf{O guloso jamais encontrará o melhor braço.} Escolhido
        $a_2$, sua média amostral converge para $0{,}9 > 0$; os outros dois
        permanecem congelados em $\Qhat = 0$ porque nunca mais são
        testados. A cirurgia --- de fato superior --- ficou eliminada por
        um único paciente azarado.
      \end{alertabox}
      \begin{ideiabox}
        A falha não é de estimação, é de \emph{cobertura}: o estimador de
        Monte Carlo é não viesado apenas para os braços que continuam
        sendo amostrados.
      \end{ideiabox}
    \end{column}
  \end{columns}
\end{frame}

%% ============================================================
\section{Arrependimento}

\begin{frame}{Um critério formal: arrependimento}
  \begin{defbox}{Arrependimento (\emph{regret})}
    Valor de ação: $\;Q(a) = \E[r \mid a]$. \quad
    Valor ótimo: $\;V^{*} = Q(a^{*}) = \max_{a \in \gA} Q(a)$.
    \smallskip

    \textbf{Perda de oportunidade} em um passo:
    \[ \ell_t = \E\bigl[V^{*} - Q(a_t)\bigr], \]
    onde a esperança é sobre a política de decisão usada para escolher $a_t$.
    \smallskip

    \textbf{Arrependimento total} até $t$:
    \[
      \Reg_t \;=\; \E\Bigl[\sum_{\tau=1}^{t} V^{*} - Q(a_\tau)\Bigr].
    \]
  \end{defbox}

  \begin{ideiabox}
    Maximizar recompensa acumulada $\iff$ minimizar arrependimento total:
    mesmo objetivo, sinal trocado, deslocamento $tV^{*}$. Mas o
    arrependimento é \emph{comparável entre problemas}.
  \end{ideiabox}
\end{frame}

%% ============================================================
\begin{frame}{Decomposição: lacunas $\times$ contagens}
  \begin{itemize}
    \item $\Nc_t(a)$: número de vezes que a ação $a$ foi escolhida até $t$.
    \item \textbf{Lacuna} $\gap_a$: diferença de valor entre $a$ e a ação
          ótima, $\;\gap_a = V^{*} - Q(a)$.
  \end{itemize}

  \[
    \Reg_t \;=\; \E\Bigl[\sum_{\tau=1}^{t} V^{*} - Q(a_\tau)\Bigr]
      \;=\; \sum_{a \in \gA} \E[\Nc_t(a)]\,\bigl(V^{*} - Q(a)\bigr)
      \;=\; \mdestaque{rlLaranja}{\sum_{a \in \gA} \E[\Nc_t(a)]\,\gap_a}
  \]

  \begin{columns}[T]
    \begin{column}{0.52\textwidth}
      \begin{ideiabox}
        Um bom algoritmo mantém \textbf{contagens pequenas onde as lacunas
        são grandes}. O problema: as lacunas são exatamente o que não
        sabemos --- se soubéssemos $\gap_a$, não haveria nada a aprender.
      \end{ideiabox}
    \end{column}
    \begin{column}{0.46\textwidth}
      \begin{alertabox}
        Toda a dificuldade da exploração está nesta frase. Não é preciso
        estimar $Q$ bem em todo braço: basta estimar bem o suficiente para
        \emph{ordenar} --- e ordenar exige menos amostras onde a lacuna é
        grande, mais onde ela é pequena.
      \end{alertabox}
    \end{column}
  \end{columns}
\end{frame}

%% ============================================================
\begin{frame}{Arrependimento do guloso, passo a passo}
  \framesubtitle{$\theta = (0{,}95;\ 0{,}90;\ 0{,}10)$, portanto $\gap = (0;\ 0{,}05;\ 0{,}85)$}

  \centering
  \begin{tabular}{ccccc}
    \toprule
    \textbf{Ação} & \textbf{Ação ótima} & \textbf{Recompensa obs.}
      & \textbf{Arrepend. do passo} & \textbf{Acumulado} \\
    \midrule
    $a_1$ & $a_1$ & 0 & 0{,}00 & 0{,}00 \\
    $a_2$ & $a_1$ & 1 & 0{,}05 & 0{,}05 \\
    $a_3$ & $a_1$ & 0 & 0{,}85 & 0{,}90 \\
    $a_2$ & $a_1$ & 1 & 0{,}05 & 0{,}95 \\
    $a_2$ & $a_1$ & 0 & 0{,}05 & 1{,}00 \\
    \bottomrule
  \end{tabular}

  \vspace{0.8em}
  \begin{columns}[T]
    \begin{column}{0.48\textwidth}
      \begin{alertabox}
        Note as linhas 2 e 4: \textbf{a recompensa observada foi $+1$ e
        mesmo assim houve arrependimento}. O arrependimento não compara com
        o que aconteceu, mas com o que aconteceria \emph{em média} sob a
        melhor ação.
      \end{alertabox}
    \end{column}
    \begin{column}{0.48\textwidth}
      \begin{ideiabox}
        Se o guloso permanecer em $a_2$, cada passo adicional custa
        $0{,}05$. Ao fim de $T$ decisões, $\Reg_T \approx 0{,}05\,T$:
        \textbf{arrependimento linear}. A curva nunca achata.
      \end{ideiabox}
    \end{column}
  \end{columns}
\end{frame}

%% ============================================================
\begin{frame}{Uma limitação incômoda --- e a saída}
  \begin{alertabox}
    Em qualquer aplicação real \textbf{não é possível calcular o
    arrependimento}: ele exige conhecer a recompensa esperada da melhor
    ação, que é justamente a incógnita. A tabela anterior só pôde ser
    preenchida porque \emph{nós} fixamos $\theta$.
  \end{alertabox}

  \begin{defbox}{O que se faz em vez disso}
    Demonstra-se uma \textbf{cota superior} para o arrependimento potencial
    de um algoritmo \emph{em qualquer problema de bandido} de uma dada
    classe. A garantia é uniforme sobre os ambientes, não estimada a
    partir dos dados.
  \end{defbox}

  \begin{columns}[T]
    \begin{column}{0.48\textwidth}
      \begin{defbox}{Cota independente do problema}
        Limita $\Reg_T$ como função apenas de $T$ e $m$
        (e.g.\ $O(\sqrt{mT\log T})$). É a garantia de \emph{pior caso}.
      \end{defbox}
    \end{column}
    \begin{column}{0.48\textwidth}
      \begin{defbox}{Cota dependente do problema}
        Limita $\Reg_T$ em função das lacunas $\gap_a$
        (e.g.\ $O(\sum_a \log T / \gap_a)$). É afiada em problemas fáceis
        e explode quando as lacunas somem.
      \end{defbox}
    \end{column}
  \end{columns}
\end{frame}

%% ============================================================
\begin{frame}{As duas cotas descrevem o mesmo algoritmo}
  \centering
  \begin{tikzpicture}
    \begin{axis}[width=0.72\textwidth, height=4.6cm,
      xlabel={$T$ (número de decisões)}, ylabel={$\Reg_T$},
      xmin=1, xmax=1000, ymin=0, ymax=60,
      axis lines=left, tick style={rlCinza}, label style={font=\scriptsize, color=rlCinza},
      tick label style={font=\tiny, color=rlCinza},
      legend style={font=\scriptsize, draw=none, fill=none, at={(0.02,0.98)}, anchor=north west},
      grid=major, grid style={rlAzul!10}]
      \addplot[rlVermelho, very thick, domain=1:1000, samples=100] {0.05*x};
      \addlegendentry{linear --- guloso, $\epsilon$ fixo}
      \addplot[rlAmbar, very thick, domain=1:1000, samples=100] {1.2*sqrt(x)};
      \addlegendentry{$\sqrt{T}$ --- pior caso do UCB}
      \addplot[rlVerde, very thick, domain=1:1000, samples=100] {4*ln(x)};
      \addlegendentry{$\log T$ --- UCB com lacunas fixas}
    \end{axis}
  \end{tikzpicture}

  \begin{ideiabox}
    ``Bom'' significa \textbf{sublinear}: $\Reg_T / T \to 0$, ou seja, a
    fração de decisões erradas tende a zero. Explorar para sempre dá
    arrependimento linear; nunca explorar também. A pergunta da aula é se
    existe algo no meio --- e a resposta é sim.
  \end{ideiabox}
\end{frame}

%% ============================================================
\section{$\epsilon$-guloso}

\begin{frame}{O algoritmo $\epsilon$-guloso}
  \begin{algbox}{$\epsilon$-guloso}
    A cada passo:
    \begin{itemize}\setlength{\itemsep}{0.2ex}
      \item com probabilidade $1 - \epsilon$, escolha
            $a_t = \argmax_{a \in \gA} \Qhat_t(a)$;
      \item com probabilidade $\epsilon$, escolha uma ação uniformemente
            ao acaso.
    \end{itemize}
  \end{algbox}

  \begin{columns}[T]
    \begin{column}{0.52\textwidth}
      \begin{itemize}
        \item Garante que \emph{todo} braço continua sendo amostrado:
              $\Nc_t(a) \to \infty$ para todo $a$.
        \item Logo $\Qhat_t(a) \to Q(a)$ para todo $a$ (lei dos grandes
              números) --- a falha de cobertura do guloso desaparece.
        \item Já usamos isto nas Aulas 4 e 6, para MDPs.
      \end{itemize}
    \end{column}
    \begin{column}{0.46\textwidth}
      \begin{alertabox}
        O preço é permanente: com $\epsilon$ fixo, o algoritmo toma uma
        decisão deliberadamente subótima em uma fração $\epsilon$ dos
        passos --- \textbf{para sempre}, muito depois de já saber qual é o
        melhor braço.
      \end{alertabox}
    \end{column}
  \end{columns}
\end{frame}

%% ============================================================
\begin{frame}{$\epsilon$-guloso no exemplo dos dedos}
  \begin{columns}[T]
    \begin{column}{0.52\textwidth}
      \begin{algbox}{Execução, com $\epsilon = 0{,}1$}
        \textbf{1.} Amostre cada braço uma vez:
        \begin{itemize}\setlength{\itemsep}{0.1ex}\small
          \item $a_1$, $r \sim \Ber(0{,}95)$: $+1$, $\Qhat(a_1) = 1$
          \item $a_2$, $r \sim \Ber(0{,}90)$: $+1$, $\Qhat(a_2) = 1$
          \item $a_3$, $r \sim \Ber(0{,}10)$: $\;\;0$, $\Qhat(a_3) = 0$
        \end{itemize}
        \textbf{2.} Qual a probabilidade de puxar cada braço em seguida?
        (empates uniformes)
      \end{algbox}
    \end{column}
    \begin{column}{0.46\textwidth}
      \pause
      \begin{respbox}
        Empate entre $a_1$ e $a_2$ na parte gulosa:
        \[
          \Prob(a_1) = \Prob(a_2)
            = \underbrace{\tfrac{0{,}9}{2}}_{\text{gulosa}}
            + \underbrace{\tfrac{0{,}1}{3}}_{\text{aleat.}}
            = 0{,}4833,
        \]
        \[
          \Prob(a_3) = \tfrac{0{,}1}{3} = 0{,}0333.
        \]
      \end{respbox}
      \pause
      \begin{ideiabox}
        \textbf{Sim, $a_3$ voltará a ser escolhido} --- e não poucas vezes:
        cerca de $T/30$ pacientes receberão ``não fazer nada'', um
        tratamento que já sabemos ser péssimo.
      \end{ideiabox}
    \end{column}
  \end{columns}
\end{frame}

%% ============================================================
\begin{frame}{Teste sua compreensão: arrependimento do $\epsilon$-guloso}
  \begin{quizbox}{}
    Lembre que $\Reg_t = \sum_{a} \E[\Nc_t(a)]\,\gap_a$, e que
    informalmente um algoritmo tem \textbf{arrependimento linear} se toma
    uma ação não ótima numa fração constante dos passos. Suponha que
    exista $a$ com $\gap_a > 0$. Quais são verdadeiras?
    \begin{enumerate}\setlength{\itemsep}{0.15ex}
      \item $\epsilon$-guloso com $\epsilon = 0{,}1$ pode ter
            arrependimento linear.
      \item $\epsilon$-guloso com $\epsilon = 0$ pode ter arrependimento
            linear.
    \end{enumerate}
    \paradiscussao
  \end{quizbox}

  \pause

  \begin{respbox}
    \textbf{As duas} --- e por motivos opostos, o que é o ponto da pergunta.
    Com $\epsilon = 0{,}1$ o algoritmo \emph{explora demais}: puxa braços
    ruins numa fração fixa $\epsilon/m$ dos passos, para sempre, e
    $\Reg_T \ge \frac{\epsilon}{m}\bigl(\sum_a \gap_a\bigr) T$. Com
    $\epsilon = 0$ ele \emph{explora de menos}: é o guloso, que pode travar
    num braço subótimo e pagar $\gap\,T$. No exemplo dos dedos, a taxa com
    $\epsilon = 0{,}1$ é $\frac{0{,}1}{3}(0{,}05 + 0{,}85) = 0{,}03$ por
    passo.
  \end{respbox}
\end{frame}

%% ============================================================
\begin{frame}{Explorar sempre e nunca explorar falham igual}
  \centering
  \begin{tikzpicture}
    \node[caixatit, text width=32mm, draw=rlVermelho!45, fill=rlVermelho!5] (n) {\textbf{Nunca explorar}\\[0.3em]
      \scriptsize $\epsilon = 0$\\[0.2em] \scriptsize trava em braço subótimo\\[0.3em]
      \scriptsize $\Reg_T = \Theta(T)$};
    \node[caixatit, text width=36mm, draw=rlVerde!55, fill=rlVerde!6, right=10mm of n] (m) {\textbf{Explorar decrescendo}\\[0.3em]
      \scriptsize $\epsilon_t \downarrow 0$ na taxa certa,\\ \scriptsize ou otimismo\\[0.3em]
      \scriptsize $\Reg_T = O(\log T)$};
    \node[caixatit, text width=32mm, draw=rlVermelho!45, fill=rlVermelho!5, right=10mm of m] (s) {\textbf{Explorar sempre}\\[0.3em]
      \scriptsize $\epsilon$ fixo\\[0.2em] \scriptsize erra fração fixa\\[0.3em]
      \scriptsize $\Reg_T = \Theta(T)$};
  \end{tikzpicture}

  \vspace{0.9em}
  \begin{teobox}{$\epsilon$ decrescente (Auer, Cesa-Bianchi e Fischer, 2002)}
    Se $\gap = \min_{a : \gap_a > 0} \gap_a$ e escolhemos
    $\displaystyle \epsilon_t = \min\Bigl\{1, \frac{c\,m}{\gap^2\,t}\Bigr\}$
    com $c > 0$ apropriado, então o $\epsilon_t$-guloso atinge
    arrependimento \textbf{logarítmico} em $T$.
  \end{teobox}

  \begin{alertabox}
    O agendamento ótimo depende de $\gap$, que é desconhecido. É por isso
    que o $\epsilon$-guloso decrescente é teoricamente correto mas
    praticamente frágil --- e por que vale a pena procurar algo que se
    adapte sozinho.
  \end{alertabox}
\end{frame}

%% ============================================================
\begin{frame}{Interlúdio: explorar-depois-comprometer}
  \framesubtitle{O algoritmo mais simples com arrependimento sublinear}

  \begin{algbox}{ETC (\emph{explore-then-commit})}
    Puxe cada braço exatamente $k$ vezes ($mk$ passos de exploração);
    depois disso, escolha para sempre o braço de maior média empírica.
  \end{algbox}

  \begin{columns}[T]
    \begin{column}{0.55\textwidth}
      Com recompensas $1$-sub-gaussianas e dois braços:
      \[
        \Reg_T \;\le\; k\gap
          \;+\; (T - 2k)\,\gap \exp\Bigl(-\frac{k\gap^2}{4}\Bigr).
      \]
      \begin{itemize}\small
        \item 1º termo: o custo pago de explorar.
        \item 2º: o risco de se comprometer com o braço errado.
        \item Otimizando $k$ e o pior $\gap$: $\Reg_T = O(T^{2/3})$.
      \end{itemize}
    \end{column}
    \begin{column}{0.43\textwidth}
      \begin{ideiabox}[fontupper=\footnotesize]
        Sublinear, portanto ``bom'' --- mas $T^{2/3}$ está longe de
        $\sqrt{T}$. A perda vem de \textbf{separar} exploração e explotação
        em fases estanques: o ETC não usa o que aprendeu \emph{durante} a
        exploração para abandonar cedo um braço ruim.
      \end{ideiabox}
      \begin{alertabox}[fontupper=\footnotesize]
        \textbf{Exploração e explotação devem ser entrelaçadas}, não
        sequenciadas.
      \end{alertabox}
    \end{column}
  \end{columns}
\end{frame}

%% ============================================================
\begin{frame}{Quão bom é possível ser? Limite inferior}
  \begin{itemize}
    \item O desempenho de \emph{qualquer} algoritmo é determinado pela
          semelhança entre o braço ótimo e os demais.
    \item Problemas difíceis têm braços que \emph{parecem} iguais mas têm
          médias diferentes: é preciso muita amostra para separá-los.
    \item Formalmente, isto é capturado pela lacuna $\gap_a$ e pela
          divergência $\KL(\mathcal{R}_a \,\|\, \mathcal{R}_{a^{*}})$.
  \end{itemize}

  \begin{teobox}{Teorema (Lai e Robbins, 1985)}
    O arrependimento total assintótico de qualquer algoritmo consistente é
    ao menos logarítmico no número de passos:
    \[
      \lim_{t \to \infty} \Reg_t \;\ge\; \log t
        \sum_{a \,:\, \gap_a > 0} \frac{\gap_a}{\KL(\mathcal{R}_a \,\|\, \mathcal{R}_{a^{*}})}.
    \]
  \end{teobox}

  \begin{ideiabox}
    Boa notícia disfarçada de má: o limite inferior é \textbf{sublinear}.
    Nada proíbe um algoritmo de acertar quase sempre --- e $\log t$ é o
    melhor que se pode esperar.
  \end{ideiabox}
\end{frame}

%% ============================================================
\begin{frame}{Lendo o limite inferior nos números do exemplo}
  Para braços de Bernoulli,
  $\KL(\theta_a \| \theta_{*}) = \theta_a \log \frac{\theta_a}{\theta_{*}}
   + (1-\theta_a)\log\frac{1-\theta_a}{1-\theta_{*}}$.

  \vspace{0.4em}
  \centering
  \begin{tabular}{lcccc}
    \toprule
    \textbf{Braço} & $\theta_a$ & $\gap_a$
      & $\KL(\theta_a\|\,0{,}95)$ & $\gap_a / \KL$ \\
    \midrule
    esparadrapo $a_2$ & 0{,}90 & 0{,}05 & 0{,}0207 & \textbf{2{,}42} \\
    nada $a_3$        & 0{,}10 & 0{,}85 & 2{,}3762 & 0{,}36 \\
    \bottomrule
  \end{tabular}

  \vspace{0.8em}
  \begin{columns}[T]
    \begin{column}{0.5\textwidth}
      \begin{alertabox}
        A inversão é o ponto: o braço \textbf{quase tão bom} ($a_2$, lacuna
        minúscula) contribui \emph{sete vezes mais} para o limite inferior
        do que o braço obviamente ruim. Distinguir $0{,}90$ de $0{,}95$
        custa muitas amostras; descartar $0{,}10$ custa poucas.
      \end{alertabox}
    \end{column}
    \begin{column}{0.48\textwidth}
      \begin{ideiabox}
        Em $T = 1000$, o piso de Lai--Robbins vale
        $(2{,}42 + 0{,}36)\log 1000 \approx 19$. Compare com o
        arrependimento do guloso, $\approx 50$, e com o do
        $\epsilon$-guloso fixo, $\approx 30$. Há espaço.
      \end{ideiabox}
    \end{column}
  \end{columns}
\end{frame}

%% ============================================================
\section{Otimismo diante da incerteza}

\begin{frame}{O princípio do otimismo}
  \begin{defbox}{Otimismo diante da incerteza}
    Escolha ações que \textbf{poderiam} ter valor alto --- não as que
    \emph{parecem} ter valor alto.
  \end{defbox}

  \pause
  Por quê? Porque há apenas dois desfechos possíveis, e ambos são bons:

  \begin{columns}[T]
    \begin{column}{0.48\textwidth}
      \begin{teobox}{Desfecho 1 --- acertamos}
        Se o braço de fato tem média alta, colhemos recompensa alta.
        Nenhum arrependimento.
      \end{teobox}
    \end{column}
    \begin{column}{0.48\textwidth}
      \begin{teobox}{Desfecho 2 --- aprendemos}
        Se o braço tem média baixa, puxá-lo reduz (em esperança) sua média
        estimada \emph{e} a incerteza sobre ela. O braço é rebaixado e não
        volta a enganar.
      \end{teobox}
    \end{column}
  \end{columns}

  \pause
  \begin{ideiabox}
    O otimismo é \textbf{autocorretivo}: só se pode estar
    excessivamente otimista sobre um braço por um número limitado de
    puxadas, porque cada puxada corrige o excesso. É exatamente esta
    contagem --- ``quantas vezes posso me enganar antes de ser
    desmentido?'' --- que se transformará na demonstração da cota.
  \end{ideiabox}
\end{frame}

%% ============================================================
\begin{frame}{Cotas superiores de confiança}
  \begin{defbox}{UCB}
    Estime uma cota superior de confiança $U_t(a)$ para cada valor de
    ação, tal que $Q(a) \le U_t(a)$ com alta probabilidade. A cota
    depende de $\Nc_t(a)$: quanto mais puxamos, mais estreita. Então
    \[
      a_t \;=\; \argmax_{a \in \gA} \; U_t(a).
    \]
  \end{defbox}

  \centering
  \begin{tikzpicture}[yscale=1.5]
    \foreach \i/\q/\w/\lab in {1/0.62/0.16/{$a_1$}, 2/0.50/0.42/{$a_2$}, 3/0.30/0.10/{$a_3$}}{
      \draw[rlAzul!30, very thick] (\i*2.3, \q-\w) -- (\i*2.3, \q+\w);
      \draw[rlAzul!30, very thick] (\i*2.3-0.12, \q-\w) -- (\i*2.3+0.12, \q-\w);
      \draw[rlAzul!30, very thick] (\i*2.3-0.12, \q+\w) -- (\i*2.3+0.12, \q+\w);
      \fill[rlAzul] (\i*2.3, \q) circle (1.6pt);
      \node[font=\scriptsize, text=rlCinza, below=1mm] at (\i*2.3, 0.05) {\lab};
      \node[font=\tiny, text=rlAzul, left=1mm] at (\i*2.3, \q) {$\Qhat$};
      \node[font=\tiny, text=rlLaranja, right=1mm] at (\i*2.3, \q+\w) {$U$};}
    \draw[rlLaranja, thick, dashed] (1.8, 0.92) -- (7.6, 0.92);
    \node[font=\scriptsize, text=rlLaranja, right] at (7.6, 0.92) {maior $U$};
    \node[font=\scriptsize, text=rlCinza, align=left, right] at (8.6, 0.5)
      {$a_2$ é escolhido:\\ estimativa \emph{menor},\\ incerteza \emph{maior}.};
  \end{tikzpicture}
\end{frame}

%% ============================================================
\begin{frame}{De onde vem a largura da barra}
  \begin{defbox}{Variável $\sigma$-sub-gaussiana}
    $X$ com $\E[X] = 0$ é $\sigma$-sub-gaussiana se
    $\E[e^{\lambda X}] \le e^{\lambda^2\sigma^2/2}$ para todo
    $\lambda \in \R$. Sua cauda decai ao menos tão rápido quanto a de uma
    gaussiana de desvio $\sigma$.
  \end{defbox}

  \begin{teobox}{Corolário 5.5 (Lattimore e Szepesvári, \emph{Bandit Algorithms})}
    Sejam $X_i - \mu$ independentes e $\sigma$-sub-gaussianas, e
    $\muhat = \frac{1}{n}\sum_{t=1}^{n} X_t$. Então, para todo
    $\varepsilon \ge 0$,
    \[
      \Prob(\muhat \ge \mu + \varepsilon) \le \exp\Bigl(-\frac{n\varepsilon^2}{2\sigma^2}\Bigr),
      \qquad
      \Prob(\muhat \le \mu - \varepsilon) \le \exp\Bigl(-\frac{n\varepsilon^2}{2\sigma^2}\Bigr).
    \]
  \end{teobox}

  \begin{ideiabox}
    Recompensas em $[0,1]$ --- como as do nosso exemplo --- são
    $\tfrac12$-sub-gaussianas pelo lema de Hoeffding. Usar $\sigma = 1$,
    como faremos, é apenas conservador: as barras saem mais largas do que
    o necessário, nunca mais estreitas.
  \end{ideiabox}
\end{frame}

%% ============================================================
\begin{frame}{Invertendo a desigualdade}
  Queremos, em vez de ``qual a probabilidade de errar por
  $\varepsilon$?'', a pergunta operacional: ``\emph{dado} um risco
  $\delta$ que aceito correr, qual $\varepsilon$?''

  \[
    \delta \;=\; \exp\Bigl(-\frac{n\varepsilon^2}{2\sigma^2}\Bigr)
    \quad\Longleftrightarrow\quad
    \varepsilon \;=\; \sqrt{\frac{2\sigma^2 \log(1/\delta)}{n}}.
  \]

  \begin{teobox}{Cota superior de confiança}
    Para qualquer $\delta \in (0,1]$, com probabilidade ao menos
    $1 - \delta$,
    \[
      \mu \;\le\; \mdestaque{rlLaranja}{\muhat + \sqrt{\frac{2\sigma^2\log(1/\delta)}{n}}}.
    \]
  \end{teobox}

  \begin{alertabox}
    Cota \emph{unilateral}. Para controlar $\abs{\muhat - \mu}$ é preciso
    aplicar as duas caudas com união de eventos: tomando
    $\delta' = \delta/2$, com probabilidade ao menos $1 - \delta$,
    $\;\abs{\mu - \muhat} \le \sqrt{2\sigma^2\log(2/\delta)/n}$.
  \end{alertabox}
\end{frame}

%% ============================================================
\begin{frame}{O algoritmo UCB1}
  \begin{algbox}{UCB1 (Auer, Cesa-Bianchi e Fischer, 2002)}
    Supondo recompensas $1$-sub-gaussianas ($\sigma^2 = 1$):
    \[
      a_t \;=\; \argmax_{a \in \gA}
        \Bigl[\; \underbrace{\Qhat(a)}_{\text{explotação}}
        \;+\; \underbrace{\sqrt{\frac{2\log(1/\delta)}{\Nc_t(a)}}}_{\text{bônus de exploração}} \;\Bigr].
    \]
  \end{algbox}

  \begin{columns}[T]
    \begin{column}{0.52\textwidth}
      \begin{itemize}
        \item Sem parâmetro de exploração para ajustar: o bônus é
              \emph{deduzido}, não escolhido.
        \item O bônus decai como $1/\sqrt{\Nc}$ --- braços pouco testados
              recebem crédito automaticamente.
        \item Braços nunca testados têm bônus infinito: o algoritmo começa
              amostrando cada braço uma vez.
      \end{itemize}
    \end{column}
    \begin{column}{0.46\textwidth}
      \begin{ideiabox}
        \textbf{A escolha de $\delta$.} Fixo, o bônus não encolhe com o
        tempo. Tomando $\delta = 1/t^2$ obtém-se a versão \emph{anytime},
        $\Qhat(a) + \sqrt{4\log t / \Nc_t(a)}$ --- a fórmula clássica com
        $\log t$ no numerador, que é a que se usa na prática.
      \end{ideiabox}
    \end{column}
  \end{columns}
\end{frame}

%% ============================================================
\begin{frame}{UCB1 no exemplo dos dedos}
  \framesubtitle{$\delta = 0{,}1$, logo o bônus vale $2{,}146/\sqrt{\Nc(a)}$; início com cada braço puxado uma vez, obtendo $(+1, +1, 0)$}

  \centering\scriptsize
  \begin{tabular}{ccccccc}
    \toprule
    $t$ & $U(a_1)$ & $U(a_2)$ & $U(a_3)$ & \textbf{escolhe} & $r$
      & \textbf{arrepend. acum.} \\
    \midrule
    3 & \textbf{3,146} & 3,146 & 2,146 & $a_1$ & $+1$ & 0,90 \\
    4 & 2,517 & \textbf{3,146} & 2,146 & $a_2$ & $+1$ & 0,95 \\
    5 & \textbf{2,517} & 2,517 & 2,146 & $a_1$ & $+1$ & 0,95 \\
    6 & 2,239 & \textbf{2,517} & 2,146 & $a_2$ & $\;\;0$ & 1,00 \\
    7 & \textbf{2,239} & 1,906 & 2,146 & $a_1$ & $+1$ & 1,00 \\
    8 & 2,073 & 1,906 & \textbf{2,146} & $a_3$ & $\;\;0$ & 1,85 \\
    9 & \textbf{2,073} & 1,906 & 1,517 & $a_1$ & $+1$ & 1,85 \\
    \bottomrule
  \end{tabular}

  \vspace{0.7em}
  \begin{columns}[T]
    \begin{column}{0.5\textwidth}
      \begin{ideiabox}
        Em $t = 8$, $a_3$ é escolhido \emph{apesar} de $\Qhat(a_3) = 0$: seu
        bônus, intocado desde a primeira puxada, ultrapassou os dos outros
        dois, que encolheram. O otimismo cobra a dívida de exploração ---
        uma vez.
      \end{ideiabox}
    \end{column}
    \begin{column}{0.48\textwidth}
      \begin{alertabox}
        Em $t = 9$, com resultado $0$, $U(a_3)$ despenca para $1{,}517$ e o
        braço sai de cena --- ao passo que o $\epsilon$-guloso voltaria a
        ele em $3{,}3\%$ dos passos \emph{indefinidamente}.
      \end{alertabox}
    \end{column}
  \end{columns}
\end{frame}

%% ============================================================
\section{A cota de arrependimento do UCB}

\begin{frame}{O nível de confiança $\delta$ é sutil}
  \begin{itemize}
    \item A cota $\Prob(\mu > \muhat + \varepsilon) \le \delta$ vale para
          \textbf{um} braço, em \textbf{um} instante, com $n$ \textbf{fixo}.
    \item O algoritmo consulta $m$ braços a cada um de $T$ passos, e $\Nc$
          é aleatório, escolhido pelo próprio algoritmo.
    \item Precisamos que \emph{nenhuma} dessas cotas falhe. União de
          eventos:
          $\;\Prob\bigl(\bigcup_i E_i\bigr) \le \sum_i \Prob(E_i)$.
  \end{itemize}

  \begin{alertabox}
    Se quisermos que todas as cotas valham simultaneamente com
    probabilidade $1 - \delta$, cada uma precisa de risco $\delta/(mT)$.
    Como o $\delta$ entra como $\log(1/\delta)$, o preço é apenas
    $\log(mT)$ --- \textbf{logarítmico no número de cotas}. É por isso que
    a união de eventos, brutal como parece, é aceitável aqui.
  \end{alertabox}

  \begin{ideiabox}
    Este é o motivo técnico da escolha $\delta = 1/n^2$ na demonstração a
    seguir: $n$ passos $\times$ risco $1/n^2$ dá risco total $1/n$, que
    contribui com $n \cdot \frac1n = 1$ puxada extra --- desprezível.
  \end{ideiabox}
\end{frame}

%% ============================================================
\begin{frame}{Cota de arrependimento do UCB --- montagem}
  Sejam $\gap_i = \mu^{*} - \mu_i$ e $\Nc_i(n)$ o número de puxadas do
  braço $i$ até o passo $n$. Já sabemos que
  \[
    \Reg_n = \sum_{i \,:\, \gap_i > 0} \gap_i \, \E[\Nc_i(n)].
  \]

  \begin{defbox}{O ``bom evento'' $G_i$}
    A média do braço $i$ nunca ultrapassa sua cota, e a cota do braço ótimo
    nunca cai abaixo de $\mu^{*}$:
    \[
      G_i \;\defeq\; \Bigl\{ \mu_i \le \min_{t} \UCB_i(t, \delta) \Bigr\}
        \cap \Bigl\{ \mu^{*} < \min_t \UCB_{a^{*}}(t, \delta) \Bigr\},
      \quad\text{com } \UCB_i(t,\delta) = \Qhat_i(t-1) + \sqrt{\tfrac{2\log(1/\delta)}{\Nc_{t-1}(i)}}.
    \]
  \end{defbox}

  \vspace{-0.4em}
  \[
    \E[\Nc_i(n)] \;\le\; \underbrace{u_i}_{\text{no bom evento}}
      \;+\; \underbrace{\Prob(G_i^{c})\, n}_{\text{fora dele}},
    \qquad u_i \defeq \frac{2\log(1/\delta)}{\gap_i^2}.
  \]
  \vspace{-0.5em}
  \begin{ideiabox}
    \textbf{No bom evento limitamos as puxadas por um argumento
    determinístico; fora dele limitamos por $n$, e mostramos que sua
    probabilidade é minúscula.}
  \end{ideiabox}
\end{frame}

%% ============================================================
\begin{frame}{Cota de arrependimento do UCB --- o passo central}
  \begin{teobox}{Afirmação}
    Se $G_i$ vale, então o braço subótimo $i \ne a^{*}$ é puxado no máximo
    $u_i = 2\log(1/\delta)/\gap_i^2$ vezes.
  \end{teobox}

  {\small
  \textbf{Demonstração (por contradição).} Suponha $G_i$ e que
  $\Nc_t(i) > u_i$ em algum instante $t$ em que $i$ seria escolhido. Como o
  bônus decresce em $\Nc$,
  \[
    \UCB_i(t,\delta) \;\le\; \mu_i + \sqrt{\frac{2\log(1/\delta)}{u_i}}
      \;=\; \mu_i + \sqrt{\gap_i^2}
      \;=\; \mu_i + \gap_i \;=\; \mu^{*}.
  \]
  Mas, por definição de $G_i$, $\;\mu^{*} < \UCB_{a^{*}}(t,\delta)$. Logo
  $\UCB_i(t,\delta) < \UCB_{a^{*}}(t,\delta)$ e o braço $i$
  \textbf{não pode} ser selecionado em $t$ --- contradição. \hfill $\square$
  }

  \begin{ideiabox}
    A conta em uma frase: \textbf{$u_i$ puxadas encolhem o bônus até ele
    caber exatamente dentro da lacuna}. A partir daí o otimismo já não
    consegue disfarçar a inferioridade do braço. Note que $u_i \propto
    1/\gap_i^2$: braços quase tão bons quanto o ótimo são puxados muitas
    vezes --- e isso é \emph{correto}, porque errar entre eles custa pouco.
  \end{ideiabox}
\end{frame}

%% ============================================================
\begin{frame}{Cota de arrependimento do UCB --- fecho}
  Tomando $\delta = 1/n^2$: o risco por cota é $1/n^2$, e a união sobre os
  $n$ instantes dá $\Prob(G_i^{c}) \le n\delta = 1/n$. Assim
  \[
    \E[\Nc_i(n)] \;\le\; \frac{2\log(n^2)}{\gap_i^2} + n \cdot \frac{1}{n}
      \;=\; \frac{4\log n}{\gap_i^2} + 1,
  \]
  e portanto

  \begin{teobox}{Cota dependente do problema}
    \[
      \Reg_n \;\lesssim\; \sum_{i \,:\, \gap_i > 0} \frac{4\log n}{\gap_i}
        \;+\; \sum_i \gap_i .
    \]
  \end{teobox}

  \begin{columns}[T]
    \begin{column}{0.5\textwidth}
      \begin{ideiabox}
        \textbf{Logarítmico em $n$} --- a ordem do piso de Lai e Robbins. O
        UCB é ótimo em taxa; a diferença está na constante.
      \end{ideiabox}
    \end{column}
    \begin{column}{0.48\textwidth}
      \begin{alertabox}
        O segundo somatório é o custo irredutível de puxar cada braço uma
        vez. E a cota \emph{explode} quando $\gap_i \to 0$: não pode ser a
        história toda.
      \end{alertabox}
    \end{column}
  \end{columns}
\end{frame}

%% ============================================================
\begin{frame}{A cota independente do problema}
  A cota anterior é inútil quando alguma lacuna é minúscula. Solução:
  \textbf{partir os braços em dois grupos} num limiar $\gap$ à escolha.
  \vspace{-0.3em}

  \[
    \Reg_n \;\le\;
      \underbrace{n\gap}_{\substack{\text{braços com } \gap_i \le \gap:\\ \text{errar custa pouco}}}
      \;+\;
      \underbrace{\sum_{i \,:\, \gap_i > \gap} \frac{4\log n}{\gap_i}}_{\substack{\text{braços com } \gap_i > \gap:\\ \text{são descartados rápido}}}
      \;\le\; n\gap + \frac{4m\log n}{\gap}.
  \]

  \vspace{-0.5em}
  Minimizando em $\gap$: $\gap^{\star} = \sqrt{4m\log n / n}$, e

  \begin{teobox}{Cota de pior caso}
    \[
      \Reg_n \;\le\; 4\sqrt{m\,n\log n} \;=\; O\bigl(\sqrt{m\,n\log n}\bigr).
    \]
  \end{teobox}

  \begin{ideiabox}
    O piso minimax é $\Omega(\sqrt{m n})$: o UCB1 fica a um fator
    $\sqrt{\log n}$ do ótimo, eliminável por variantes (MOSS, KL-UCB). E
    \textbf{o mesmo algoritmo} satisfaz as duas cotas: adapta-se sozinho à
    dificuldade do problema.
  \end{ideiabox}
\end{frame}

%% ============================================================
\begin{frame}{Verificando as cotas nos números do exemplo}
  \framesubtitle{$m = 3$, $\gap = (0;\ 0{,}05;\ 0{,}85)$, $n = 1000$}

  \centering\small
  \begin{tabular}{lr}
    \toprule
    \textbf{Referência} & $\Reg_{1000}$ \\
    \midrule
    Piso de Lai--Robbins & $\approx 19$ \\
    Cota do UCB, dependente do problema & $\approx 586$ \\
    Cota do UCB, independente do problema & $\approx 576$ \\
    Guloso (travado em $a_2$) & $\approx 50$ \\
    $\epsilon$-guloso, $\epsilon = 0{,}1$ & $\approx 30$ \\
    \bottomrule
  \end{tabular}

  \vspace{0.8em}
  \begin{alertabox}[fontupper=\footnotesize]
    \textbf{As cotas do UCB são piores que o desempenho observado dos
    métodos ingênuos} --- e isso não é erro: é o que significa uma cota de
    \emph{pior caso}. Ela vale para todo bandido com estas lacunas; os
    números do guloso valem para \emph{esta} realização.
  \end{alertabox}
  \begin{ideiabox}[fontupper=\footnotesize]
    A comparação honesta é assintótica: em $n = 10^6$ o guloso paga
    $5\times10^4$ e o $\epsilon$-guloso $3\times10^4$, ambos lineares,
    contra $\approx 1{,}2\times10^3$ da cota do UCB. As curvas se cruzam ---
    e nunca mais se reencontram.
  \end{ideiabox}
\end{frame}

%% ============================================================
\section{Variações e encerramento}

\begin{frame}{Teste sua compreensão: e se usássemos a cota \emph{inferior}?}
  \begin{quizbox}{}
    Uma alternativa seria escolher sempre o braço de maior \textbf{cota
    inferior} de confiança,
    $\;a_t = \argmax_a \bigl[\Qhat(a) - \sqrt{2\log(1/\delta)/\Nc_t(a)}\bigr]$.
    Isso garante arrependimento baixo? Considere o caso de dois braços.
    \paradiscussao
  \end{quizbox}

  \pause

  \begin{respbox}[fontupper=\footnotesize]
    \textbf{Não --- é pior que o guloso.} O bônus vira \emph{penalidade}
    por incerteza: braços pouco testados ficam artificialmente ruins, e são
    justamente os que precisavam ser testados. Com
    $\mu_1 = 0{,}5$, $\mu_2 = 0{,}6$ e puxadas iniciais $r_1 = 1$,
    $r_2 = 0$, temos $\LCB_1 = 1 - c > \LCB_2 = -c$. Cada nova puxada de
    $a_1$ encolhe sua penalidade ($\LCB_1 \to 0{,}5$), enquanto $\LCB_2$
    fica congelado. O braço $2$ nunca mais é testado: arrependimento
    $0{,}1\,T$.
  \end{respbox}

  \begin{ideiabox}[fontupper=\footnotesize]
    \textbf{O otimismo é autocorretivo; o pessimismo, autoconfirmatório.}
    (O pessimismo tem seu lugar em RL \emph{offline}, onde não se pode
    testar hipóteses e evitar o desconhecido é o que se deseja.)
  \end{ideiabox}
\end{frame}

%% ============================================================
\begin{frame}{Uma segunda família: casamento de probabilidade}
  \begin{defbox}{Arrependimento bayesiano}
    Coloque uma \emph{priori} $p(\theta)$ sobre os parâmetros do bandido e
    meça a média do arrependimento sob ela:
    \[
      \mathrm{BayesReg}_T \;=\; \E_{\theta \sim p}\bigl[\Reg_T(\theta)\bigr].
    \]
    Critério mais brando que o de pior caso --- e que admite algoritmos
    notavelmente simples.
  \end{defbox}

  \begin{algbox}{Amostragem de Thompson (Thompson, 1933)}
    A cada passo: \textbf{(1)} amostre um parâmetro da posteriori,
    $\tilde\theta_a \sim p(\theta_a \mid \text{histórico})$ para cada $a$;
    \textbf{(2)} aja \emph{como se} $\tilde\theta$ fosse a verdade,
    $a_t = \argmax_a Q_{\tilde\theta}(a)$; \textbf{(3)} observe $r_t$ e
    atualize a posteriori.
  \end{algbox}

  \begin{ideiabox}
    O nome vem da propriedade que a define: cada braço é escolhido com
    probabilidade \emph{igual à probabilidade posterior de ele ser o
    ótimo}. A exploração não é imposta por um bônus --- ela emerge da
    incerteza da posteriori, e desaparece sozinha quando a incerteza some.
  \end{ideiabox}
\end{frame}

%% ============================================================
\begin{frame}{Thompson com braços de Bernoulli}
  Priori conjugada: $\theta_a \sim \Bet(\alpha_a, \beta_a)$. Após observar
  $r$, a atualização é contagem pura:
  \[
    \alpha_a \leftarrow \alpha_a + r, \qquad \beta_a \leftarrow \beta_a + (1 - r).
  \]

  \begin{columns}[T]
    \begin{column}{0.5\textwidth}
      Com priori uniforme $\Bet(1,1)$ e as mesmas três primeiras puxadas
      $(+1, +1, 0)$:
      \smallskip

      \centering\small
      \begin{tabular}{lccc}
        \toprule
        \textbf{Braço} & \textbf{Posteriori} & $\E[\theta]$
          & $\Prob(\text{é o melhor})$ \\
        \midrule
        $a_1$ & $\Bet(2,1)$ & 0,667 & 0,467 \\
        $a_2$ & $\Bet(2,1)$ & 0,667 & 0,467 \\
        $a_3$ & $\Bet(1,2)$ & 0,333 & \textbf{0,066} \\
        \bottomrule
      \end{tabular}
    \end{column}
    \begin{column}{0.48\textwidth}
      \begin{ideiabox}
        Compare com o $\epsilon$-guloso, que daria $3{,}3\%$ a $a_3$
        \emph{para sempre}. Thompson dá $6{,}6\%$ \emph{agora}, quando ainda
        há dúvida legítima, e essa fração cai sozinha à medida que a
        posteriori de $a_3$ se concentra perto de $0{,}1$.
      \end{ideiabox}
      \begin{alertabox}
        Note que $\Prob(\text{melhor}) \ne$ ordenação das médias: $a_3$ tem
        média posterior bem menor, mas $6{,}6\%$ de chance de ser o melhor
        --- com uma única observação, a cauda ainda é gorda.
      \end{alertabox}
    \end{column}
  \end{columns}

\end{frame}

%% ------------------------------------------------------------
\begin{frame}{Garantias e comparação com o otimismo}
  \begin{teobox}{Garantias da amostragem de Thompson}
    Arrependimento bayesiano $O(\sqrt{mT\log T})$. E, embora bayesiano por
    construção, Thompson também satisfaz cotas \emph{frequentistas}
    $O(\sum_i \log T / \gap_i)$, sendo assintoticamente ótimo para braços
    de Bernoulli (Agrawal e Goyal, 2012; Kaufmann \emph{et al.}, 2012).
  \end{teobox}

  \begin{columns}[T]
    \begin{column}{0.48\textwidth}
      \begin{defbox}{A simetria com o UCB}
        Ambos escolhem por um valor \emph{plausível}, não pelo esperado. O
        UCB toma o \textbf{quantil superior} da incerteza; Thompson toma
        uma \textbf{amostra} dela. Um é determinístico e pessimista quanto
        ao risco; o outro é aleatório e calibrado.
      \end{defbox}
    \end{column}
    \begin{column}{0.48\textwidth}
      \begin{alertabox}
        A aleatoriedade de Thompson é uma vantagem prática subestimada: em
        sistemas com decisões \emph{em lote} (mil usuários atendidos com a
        mesma política), o UCB daria a todos o mesmo braço, enquanto
        Thompson diversifica automaticamente.
      \end{alertabox}
    \end{column}
  \end{columns}
\end{frame}

%% ============================================================
\begin{frame}{Panorama comparativo}
  \centering\small
  \begin{tabular}{lccl}
    \toprule
    \textbf{Algoritmo} & \textbf{Arrepend.} & \textbf{Parâmetro}
      & \textbf{Por que falha ou funciona} \\
    \midrule
    Guloso & $\Theta(T)$ & --- & trava; nunca corrige um azar inicial \\
    $\epsilon$-guloso, $\epsilon$ fixo & $\Theta(T)$ & $\epsilon$
      & explora uniformemente, para sempre \\
    $\epsilon_t$-guloso, $\epsilon_t \propto 1/t$ & $O(\log T)$ & $c, \gap$
      & ótimo em taxa, mas exige conhecer $\gap$ \\
    ETC & $O(T^{2/3})$ & $k$ & separa as fases; não se adapta \\
    \textbf{UCB1} & $O(\log T)$ & $\delta$ (ou nenhum)
      & bônus deduzido; adapta-se às lacunas \\
    \textbf{Thompson} & $O(\log T)$ & priori
      & explora na medida da dúvida posterior \\
    \midrule
    \emph{Piso (Lai--Robbins)} & $\Omega(\log T)$ & ---
      & \emph{ninguém faz melhor} \\
    \bottomrule
  \end{tabular}

  \vspace{0.9em}
  \begin{ideiabox}
    Empiricamente, Thompson costuma bater o UCB1 apesar de garantias
    semelhantes --- o bônus do UCB é calibrado para o pior caso, enquanto a
    posteriori usa a forma real da incerteza. Em compensação, o UCB não
    exige escolher uma priori nem amostrar dela.
  \end{ideiabox}
\end{frame}

%% ============================================================
\begin{frame}{Resumo da aula}
  \begin{itemize}
    \item O \textbf{bandido multi-braços} isola a exploração ao remover as
          consequências adiadas: um MDP de um estado só.
    \item O \textbf{arrependimento} $\Reg_T = \sum_a \E[\Nc_T(a)]\gap_a$
          mede o custo da aprendizagem. ``Bom'' $=$ sublinear.
    \item O \textbf{guloso} tem arrependimento linear por falta de
          cobertura; o \textbf{$\epsilon$-guloso fixo}, por excesso
          permanente de exploração.
    \item \textbf{Lai e Robbins}: nenhum algoritmo faz melhor que
          $\Omega(\log T)$, com constante $\gap_a/\KL$.
    \item O \textbf{otimismo} resolve o dilema sem conhecer as lacunas: o
          bônus $\sqrt{2\log(1/\delta)/\Nc}$ sai da concentração, e a
          demonstração conta quantas puxadas bastam para o bônus caber
          dentro da lacuna.
    \item \textbf{UCB1}: $O(\sum_i \log T/\gap_i)$ e $O(\sqrt{mT\log T})$
          --- ótimo em taxa nos dois arcabouços.
  \end{itemize}

  \begin{ideiabox}
    Bandidos são o lugar mais simples para ver estas ideias, mas elas se
    estendem a MDPs quase sem alteração. \textbf{Próxima aula:}
    exploração em MDPs --- aprendizagem PAC, otimismo em modelos
    (\textsc{R-max}, UCRL) e cotas de complexidade amostral.
  \end{ideiabox}
\end{frame}

%% ============================================================
\begin{frame}{Créditos e leituras}
  \begin{itemize}
    \item Material adaptado e expandido de \textbf{CS234: Reinforcement
          Learning}, Emma Brunskill, Stanford University; parte da
          estrutura remonta às notas de David Silver (UCL/DeepMind).
    \item \textbf{Referência central:} Lattimore \& Szepesvári,
          \emph{Bandit Algorithms} (Cambridge, 2020) --- Capítulos 4 a 8
          cobrem toda a aula, com as demonstrações completas.
    \item Sutton \& Barto, \emph{Reinforcement Learning: An Introduction},
          2\textsuperscript{a} ed., Capítulo 2.
    \item \textbf{Artigos:} Lai \& Robbins (1985); Auer, Cesa-Bianchi \&
          Fischer, \emph{Machine Learning} (2002); Thompson (1933);
          Agrawal \& Goyal (2012); Kaufmann, Korda \& Munos (2012);
          Russo \emph{et al.}, \emph{A Tutorial on Thompson Sampling} (2018).
    \item \textbf{Expansões desta versão} em relação ao original: bandidos
          contextuais, explorar-depois-comprometer, leitura numérica do
          piso de Lai--Robbins, cota independente do problema com a
          otimização do limiar, contraexemplo detalhado da cota inferior,
          e a seção de amostragem de Thompson.
  \end{itemize}
\end{frame}

\end{document}
